% GENERATED FILE. Edit canonical lessons, then run npm run generate:books.
% canonical-source-hash: fnv1a64:d1e0ea9235d4cf18
% canonical-lessons: ES-C739-curar, ES-C739-mejorarse, ES-C739-recuperacion, ES-C739-operar, ES-C739-cirugia

\chapter{Curarse}
\label{ch:es-curarse}
\hlchaptermodality{739}

\begin{chapteropening}
\textbf{By the end of this chapter:} \emph{I can talk about treating an illness, an operation and getting better.}

Complete the last lesson of chapter 739: I can talk about treating an illness, an operation and getting better.
\end{chapteropening}

\section[curar]{curar — to cure, to heal}

\label{lesson:ES-C739-curar}



\begin{quote}
Before the new one: say the Spanish for the treatment, then the Spanish for to suffer from.
\end{quote}

\subsection*{You'll want to know: curar}
\textbf{curar} — \textquotedblleft{}to cure, to heal\textquotedblright{}.

\begin{quote}
El tiempo lo cura todo.
\end{quote}

\begin{quote}
\emph{Time heals everything.}
\end{quote}

\emph{Curarse} — to get better.

\begin{cousinweb}
\emph{curar} is inherited from Latin \emph{cūrāre} \textquotedblleft{}to take care of\textquotedblright{} (EtymDB), built on \emph{cūra} \textquotedblleft{}care\textquotedblright{}.
\end{cousinweb}

\begin{grammarlens}[title={What you've built}]
One more word for telling, explaining and describing at B1.
\end{grammarlens}

\subsection*{Your turn}
\begin{itemize}
\raggedright
  \item \emph{Say it:} \emph{curar}
  \item \emph{Say it:} \emph{curar}, once more
  \item \emph{Say it:} a sentence of your own with \emph{curar}
  \item \emph{Recall:} say the Spanish for the treatment, then the Spanish for to suffer from, then say \emph{curar} again
\end{itemize}

\begin{tcolorbox}[breakable,colback=teal!4,colframe=teal!35!black,title={Before you move on}]
What does curar mean? (\textquotedblleft{}To cure, to heal\textquotedblright{}.) Say it once more.
\end{tcolorbox}
\section[mejorarse]{mejorarse — to get better}

\label{lesson:ES-C739-mejorarse}



\begin{quote}
Before the new one: say the Spanish for to suffer from, then the Spanish for to cure.
\end{quote}

\subsection*{You'll want to know: mejorarse}
\textbf{mejorarse} — \textquotedblleft{}to get better\textquotedblright{}.

\begin{quote}
¡Que te mejores pronto!
\end{quote}

\begin{quote}
\emph{Get well soon!}
\end{quote}

\emph{¡Que te mejores!} takes the subjunctive.

\begin{cousinweb}
\emph{mejorarse} is built on \emph{mejor} \textquotedblleft{}better\textquotedblright{} (EtymDB), which is inherited from Latin \emph{melior}.
\end{cousinweb}

\begin{grammarlens}[title={What you've built}]
One more word for telling, explaining and describing at B1.
\end{grammarlens}

\subsection*{Your turn}
\begin{itemize}
\raggedright
  \item \emph{Say it:} \emph{mejorarse}
  \item \emph{Say it:} \emph{mejorarse}, once more
  \item \emph{Say it:} a sentence of your own with \emph{mejorarse}
  \item \emph{Recall:} say the Spanish for to suffer from, then the Spanish for to cure, then say \emph{mejorarse} again
\end{itemize}

\begin{tcolorbox}[breakable,colback=teal!4,colframe=teal!35!black,title={Before you move on}]
What does mejorarse mean? (\textquotedblleft{}To get better\textquotedblright{}.) Say it once more.
\end{tcolorbox}
\section[la recuperación]{la recuperación — recovery}

\label{lesson:ES-C739-recuperacion}



\begin{quote}
Before the new one: say the Spanish for to cure, then the Spanish for to get better.
\end{quote}

\subsection*{You'll want to know: la recuperación}
\textbf{la recuperación} — \textquotedblleft{}recovery\textquotedblright{}.

\begin{quote}
Le deseo una pronta recuperación.
\end{quote}

\begin{quote}
\emph{I wish you a speedy recovery.}
\end{quote}

\emph{Pronta recuperación} — get well soon.

\begin{cousinweb}
\emph{recuperación} is usually described as a learned borrowing from Latin \emph{recuperātiō}, built on \emph{recuperāre} \textquotedblleft{}to get back\textquotedblright{}.
\end{cousinweb}

\begin{grammarlens}[title={What you've built}]
One more word for telling, explaining and describing at B1.
\end{grammarlens}

\subsection*{Your turn}
\begin{itemize}
\raggedright
  \item \emph{Say it:} \emph{la recuperación}
  \item \emph{Say it:} \emph{la recuperación}, once more
  \item \emph{Say it:} a sentence of your own with \emph{la recuperación}
  \item \emph{Recall:} say the Spanish for to cure, then the Spanish for to get better, then say \emph{la recuperación} again
\end{itemize}

\begin{tcolorbox}[breakable,colback=teal!4,colframe=teal!35!black,title={Before you move on}]
What does la recuperación mean? (\textquotedblleft{}Recovery\textquotedblright{}.) Say it once more.
\end{tcolorbox}
\section[operar]{operar — to operate (on)}

\label{lesson:ES-C739-operar}



\begin{quote}
Before the new one: say the Spanish for to get better, then the Spanish for recovery.
\end{quote}

\subsection*{You'll want to know: operar}
\textbf{operar} — \textquotedblleft{}to operate (on)\textquotedblright{}.

\begin{quote}
Me operan de la rodilla en mayo.
\end{quote}

\begin{quote}
\emph{I am having a knee operation in May.}
\end{quote}

\emph{Operarse de} — to have an operation on.

\begin{cousinweb}
\emph{operar} is a learned borrowing from Late Latin \emph{operāre} \textquotedblleft{}to work\textquotedblright{} (EtymDB), built on \emph{opus} \textquotedblleft{}work\textquotedblright{}.
\end{cousinweb}

\begin{grammarlens}[title={What you've built}]
One more word for telling, explaining and describing at B1.
\end{grammarlens}

\subsection*{Your turn}
\begin{itemize}
\raggedright
  \item \emph{Say it:} \emph{operar}
  \item \emph{Say it:} \emph{operar}, once more
  \item \emph{Say it:} a sentence of your own with \emph{operar}
  \item \emph{Recall:} say the Spanish for to get better, then the Spanish for recovery, then say \emph{operar} again
\end{itemize}

\begin{tcolorbox}[breakable,colback=teal!4,colframe=teal!35!black,title={Before you move on}]
What does operar mean? (\textquotedblleft{}To operate (on)\textquotedblright{}.) Say it once more.
\end{tcolorbox}
\section[la cirugía]{la cirugía — surgery}

\label{lesson:ES-C739-cirugia}



\begin{quote}
Before the new one: say the Spanish for recovery, then the Spanish for to operate.
\end{quote}

\subsection*{You'll want to know: la cirugía}
\textbf{la cirugía} — \textquotedblleft{}surgery\textquotedblright{}.

\begin{quote}
Necesita cirugía urgente.
\end{quote}

\begin{quote}
\emph{He needs urgent surgery.}
\end{quote}

\emph{Cirugía estética} — cosmetic surgery. The surgeon is \emph{el cirujano}.

\begin{cousinweb}
\emph{cirugía} is usually traced to Latin \emph{chīrurgia}, borrowed from Greek \emph{kheirourgía} \textquotedblleft{}handwork\textquotedblright{}, built on \emph{kheír} \textquotedblleft{}hand\textquotedblright{} and \emph{érgon} \textquotedblleft{}work\textquotedblright{}.
\end{cousinweb}

\begin{grammarlens}[title={What you've built}]
One more word for telling, explaining and describing at B1.
\end{grammarlens}

\subsection*{Your turn}
\begin{itemize}
\raggedright
  \item \emph{Say it:} \emph{la cirugía}
  \item \emph{Say it:} \emph{la cirugía}, once more
  \item \emph{Say it:} a sentence of your own with \emph{la cirugía}
  \item \emph{Recall:} say the Spanish for recovery, then the Spanish for to operate, then say \emph{la cirugía} again
\end{itemize}

\begin{tcolorbox}[breakable,colback=teal!4,colframe=teal!35!black,title={Before you move on}]
What does la cirugía mean? (\textquotedblleft{}Surgery\textquotedblright{}.) Say it once more.
\end{tcolorbox}
